%%!TeX program = xelatex
%% 부득이하게 pdflatex을 사용해야 할 경우 위의 magic comment를 제거하십시오.

% Initiated by 정민석(2014년도 경기과학고 수학과전문교원)
% Continously being modified by 경기과학고 TeX 사용자협회
% Website : http://latex.gs.hs.kr
% 박기현(선생님)과 한수민(37기) 학생이 본문 나누어진 버전으로 개선함.
% 채상미(선생님)과 권준우(37기) 학생이 2021년 최종 수정.

\documentclass{gshs_thesis}

\graphicspath{{images/}}
% 이곳에 필요한 별도의 패키지들을 적어넣으시오.
%\usepackage{...}
\usepackage{verbatim} % for commment, verbatim environment
\usepackage{spverbatim} % automatic linebreak verbatim environment
\usepackage{listings}
\usepackage{indentfirst}
\lstset{
	basicstyle=\small\ttfamily,
	columns=flexible,
	breaklines=true
}
\usepackage{hologo}

% -----------------------------------------------------------------------
%                   이 부분은 수정하지 마시오.
% -----------------------------------------------------------------------
\titleheader{졸업논문청구논문}
\school{과학영재학교 경기과학고등학교}
\approval{위 논문은 과학영재학교 경기과학고등학교 졸업논문으로\\
졸업논문심사위원회에서 심사 통과하였음.}
\chairperson{심사위원장}
\examiner{심사위원}
\apprvsign{(인)}
\korabstract{초 록}
\koracknowledgement{감사의 글}
\korresearches{연 구 활 동}

% -----------------------------------------------------------------------
%                   영어로 작성할 경우 주석 처리를 제거하시오
% -----------------------------------------------------------------------
% \def\useenglish{true}

%: ----------------------------------------------------------------------
%:                  논문 제목과 저자 이름을 입력하시오
% ----------------------------------------------------------------------
\title{한글 제목} %한글 제목
\engtitle{English Title} %영문 제목
\korname{홍 길 동} %저자 이름을 한글로 입력하시오 (글자 사이 띄어쓰기)
\engname{Hong, Gil-Dong} %저자 이름을 영어로 입력하시오 (family name, personal name)
\chnname{洪 吉 東} %저자 이름을 한자로 입력하시오 (글자 사이 띄어쓰기)
\studid{21000} %학번을 입력하시오

%------------------------------------------------------------------------
%                  심사위원과 논문 승인 날짜를 입력하시오
%------------------------------------------------------------------------
\advisor{Mok, Chinook}  %지도교사 영문 이름 (family name, personal name)
\judgeone{박 승 원} %심사위원장
\judgetwo{이 주 찬}   %심사위원1
\judgethree{목 진 욱} %심사위원2(지도교사)
\degreeyear{2024}   %졸업 년도
\degreedate{2023}{6}{30} %논문 승인 날짜 양식

%------------------------------------------------------------------------
%                  논문제출 전 체크리스트를 확인하시오
%------------------------------------------------------------------------
\checklisttitle{[논문제출 전 체크리스트]} %수정하지 마시오
\checklistI{1. 이 논문은 내가 직접 연구하고 작성한 것이다.} %수정하지 마시오
% 이 항목이 사실이라면 다음 줄 앞에 "%"기호 삽입, 다다음 줄 앞의 "%"기호 제거하시오
\checklistmarkI{$\square$}
%\checklistmarkI{$\text{\rlap{$\checkmark$}}\square$}
\checklistII{2. 인용한 모든 자료(책, 논문, 인터넷자료 등)의 인용표시를 바르게 하였다.} %수정하지 마시오
% 이 항목이 사실이라면 다음 줄 앞에 "%"기호 삽입, 다다음 줄 앞의 "%"기호 제거하시오
\checklistmarkII{$\square$}
%\checklistmarkII{$\text{\rlap{$\checkmark$}}\square$}
\checklistIII{3. 인용한 자료의 표현이나 내용을 왜곡하지 않았다.} %수정하지마시오
% 이 항목이 사실이라면 다음 줄 앞에 "%"기호 삽입, 다다음 줄 앞의 "%"기호 제거하시오
\checklistmarkIII{$\square$}
%\checklistmarkIII{$\text{\rlap{$\checkmark$}}\square$}
\checklistIV{4. 정확한 출처제시 없이 다른 사람의 글이나 아이디어를 가져오지 않았다.} %수정하지 마시오
% 이 항목이 사실이라면 다음 줄 앞에 "%"기호 삽입, 다다음 줄 앞의 "%"기호 제거하시오
\checklistmarkIV{$\square$}
%\checklistmarkIV{$\text{\rlap{$\checkmark$}}\square$}
\checklistV{5. 논문 작성 중 도표나 데이터를 조작(위조 혹은 변조)하지 않았다.} %수정하지 마시오
% 이 항목이 사실이라면 다음 줄 앞에 "%"기호 삽입, 다다음 줄 앞의 "%"기호 제거하시오
\checklistmarkV{$\square$}
%\checklistmarkV{$\text{\rlap{$\checkmark$}}\square$}
\checklistVI{6. 다른 친구와 같은 내용의 논문을 제출하지 않았다.} %수정하지 마시오
% 이 항목이 사실이라면 다음 줄 앞에 "%"기호 삽입, 다다음 줄 앞의 "%"기호 제거하시오
\checklistmarkVI{$\square$}
%\checklistmarkVI{$\text{\rlap{$\checkmark$}}\square$}
 % usepackage 등의 명령어는 여기에.
\usepackage{cite}
\usepackage{textcomp}
\usepackage{tocloft}
\setlength{\cftbeforesecskip}{0pt}
\setlength{\cftbeforesubsecskip}{0pt}
\setlength{\cftbeforesubsubsecskip}{0pt}

% 본문 시작
\begin{document}

	%표지만들기
	%makecover 함수와 관련하여 "Underfull \hbox (badness 10000) in paragraph" 오류는 무시하십시오. (TeXstudio ver 2.9.4 오류 기준)
	%\makecover
	

	%\baselineskip=2.2em         % line spacing in the paragraph
	\begin{abstract}
이 부분을 수정하려면 /frontmatter/abstract를 수정하세요.~\\\\
\end{abstract}
\begin{dedication2}
{\large{19000 김경곽}}\\[5mm]
2021년 수학세미나I 0반
\end{dedication2}
\begin{copyright2}
	
	\begin{figure}[h]
		\includegraphics[width=3in]{logo_full}
	\end{figure}
	
	\noindent Copyright \textcopyright\the\year\ Junu Kwon \\ % Copyright of the form
	
	\noindent \textsc{https://www.gs.hs.kr/}\\ % URL
	
	\noindent \textit{Compiled \today} 
	\if@openright\cleardoublepage\else\clearpage\fi
\end{copyright2} % Abstract
	
	%%%%%%%%%%%%%%%%%%%%%%%%%%%%%%%%%%%%%%%%%%%%%%%%%%%%%%%%%%%
	%%%% Main Document %%%%%%%%%%%%%%%%%%%%%%%%%%%%%%%%%%%%%%%%
	%%%%%%%%%%%%%%%%%%%%%%%%%%%%%%%%%%%%%%%%%%%%%%%%%%%%%%%%%%%
	\cleardoublepage
	\clearpage
	\renewcommand{\thepage}{\arabic{page}}
	\setcounter{page}{1}
	
	%각 장을 아래와 같이 sub 폴더 안에 만들어서 넣으면 편리하다.
	\section{서론}

\subsection{BibTeX}

BibTeX을 사용하면 참고문헌을 쉽게 정리하거나 관리할 수 있다. 다음은 Google 학술검색 사이트에서 찾은 논문에서 BibTeX 코드를 찾는 방법이다.

\begin{enumerate}
	\item Google 학술검색에서 원하는 논문을 찾는다.
	\item 인용 버튼을 누르면 MLA, APA 등 다양한 인용 스타일이 나온다. 여기서 창 하단의 BibTeX을 클릭한다.
	\item BibTeX을 클릭해서 나온 코드 전문을 복사하여 보고서에 있는 폴더의 bibfile.bib 파일에 붙여넣기 한다.
\end{enumerate}

다음은 예시로 찾은 논문의 BibTeX 코드이다.
\begin{lstlisting}
	@article{ward2001landslide,
	title={Landslide tsunami},
	author={Ward, Steven N},
	journal={Journal of Geophysical Research: Solid Earth},
	volume={106},
	number={B6},
	pages={11201--11215},
	year={2001},
	publisher={Wiley Online Library}
	}
\end{lstlisting}

BibTeX에 넣은 논문을 인용하기 위해서는 cite 명령어를 사용한다. 위 코드에서 @article 오른쪽에 있는 문구가 ward2001landslide이므로 이 논문을 인용하기 위해서 명령어
\begin{lstlisting}
	\cite{ward2001landslide}
\end{lstlisting}
를 사용한다. 이 명령어 실행 시 명령어를 작성한 자리에 \cite{ward2001landslide}과 같이 논문을 인용했다는 표시가 뜬다. 인용한 논문의 목록은 문서의 References에 표시된다. % 서론
	\section{Theoretical Background}

\subsection{줄 바꿈}

\LaTeX 에서 줄 바꿈을 하는 방법은 한컴이나 워드와는 다르다.
줄 바꿈을 하기 위해서는 Enter 키를 2번 연속으로 누르거나, 
Back slash 키를 두 번 연속 누르는 방법(\textbackslash \textbackslash)을 사용한다. 
또한 여러 줄을 띄우고 싶다면, \textbackslash vskip 명령어를 사용한다. 
예를 들어, \textbackslash vskip 2pc을 사용 시
\vskip 2pc
위와 같이 한 번에 두 줄을 띄운다. (pc 대신 cm나 mm 등도 사용 가능하다)


그런데 \textbackslash \textbackslash 를 제외한 나머지 방법 사용 시 첫 단어에 들여쓰기가 되는 것을 볼 수 있다. 들여쓰기를 하지 않으려면 문단의 첫부분에 \textbackslash noindent 명령어를 사용한다.

\subsection{그림 삽입}

\LaTeX 에서 그림을 넣는 명령어는 아래와 같다.
\begin{lstlisting}
	\begin{figure}[t]
		\begin{center}
			\includegraphics[width=.3\textwidth]{figure.png}
			\caption{Figure Caption}
			\label{figlabel}
		\end{center}
	\end{figure}
\end{lstlisting}
그림을 넣기 위해서는 images 폴더에 그림이 있어야 한다.
위 명령어에서 width는 그림의 너비, caption은 그림 밑에 표시될 설명, label은 그림을 언급할 때 사용하는 그림의 코드이다. 그림의 언급은 \textbackslash ref\{figlabel\}과 같은 방법으로 한다.

\noindent center 옆의 [t]는 그림을 넣을 위치이다. h, t, b, p 총 네 가지의 설정이 가능하며 각각 here(여기), top(페이지 맨 위), bottom(페이지 맨 아래), page(새로운 페이지)를 의미한다.
일반적으로 htbp를 사용하지만 논문과 같이 페이지 상단에 그림을 위치시켜야 하는 경우 [t]를 사용한다.

\newpage

\begin{figure}[t]
	\begin{center}
		\includegraphics[width=.3\textwidth]{figure.png}
		\caption{Figure Caption}
		\label{figlabel}
	\end{center}
\end{figure}

\noindent 예시 코드 실행 시 위와 같이 페이지 상단에 그림 \ref{figlabel}이 나타난다.

\subsection{표 삽입}

\LaTeX 에서 직접 표를 만들 수 있으나 명령어가 길고 복잡하며, 표의 모습을 보면서 작업할 수 없기에 본인이 원하는 모양대로 만드는 것이 불편하다는 단점이 있다.
그러므로 표의 경우 다른 조판 프로그램에서 제작한 후 이미지 형태로 저장해 \LaTeX 문서에 넣거나, 인터넷 사이트인 Tables Generator 등을 사용하는 것이 더 효율적이다. 다음은 Tables Generator를 사용해 만든 표이다.
\begin{lstlisting}
\begin{table}[htbp]
	\centering
	\caption{렌즈의 종류}
	\label{lens}
	\resizebox{.5\textwidth}{!}{%
		\begin{tabular}{|l|l|l|}
			\hline
			\multicolumn{1}{|c|}{} & 광각렌즈 & 망원렌즈 \\ \hline
			초점거리                   & 짧다   & 길다   \\ \hline
			화각                     & 크다   & 작다   \\ \hline
		\end{tabular}%
	}
\end{table}
\end{lstlisting}

\begin{table}[htbp]
	\centering
	\caption{렌즈의 종류}
	\label{lens}
	\resizebox{.5\textwidth}{!}{%
		\begin{tabular}{|l|l|l|}
			\hline
			\multicolumn{1}{|c|}{} & 광각렌즈 & 망원렌즈 \\ \hline
			초점거리                   & 짧다   & 길다   \\ \hline
			화각                     & 크다   & 작다   \\ \hline
		\end{tabular}%
	}
\end{table}

\subsection{수식 삽입}

문장 내부에서 수식을 사용하기 위해서는 \$ 수식 \$이나 \textbackslash( 수식 \textbackslash)와 같은 방법을 사용한다.
예를 들어 '소비 전력 \$P=I\^{}2 R\$이다.'라는 코드를 작성 후 컴파일하면
\begin{quote}
'소비 전력 $P=I^2 R$이다.'
\end{quote}
\noindent 와 같이 수식이 나타난다. \\
이때 수식과 문장의 높이를 맞추기 위해 분수나 시그마 등의 수학 기호의 크기가 조절된다. 수식의 가독성을 키우기 위해서는 수식 앞에 \textbackslash displaystyle을 붙이면 된다.
\begin{itemize}
	\item displaystyle을 사용하지 않은 경우 : $\sum_{i=1}^{n} = \frac{n(n+1)}{2}$
	\item displaystyle을 사용한 경우 : \(\displaystyle \sum_{i=1}^{n} = \frac{n(n+1)}{2}\)
\end{itemize}

\vskip 2pc

\noindent 수식이 문장 안에 있는 것이 아니라 한 줄을 통째로 차지하는 경우 \textbackslash[ 수식 \textbackslash]이나 명령어 \textbackslash equation을 사용한다. 차이점이라면 전자는 수식에 번호가 붙지 않고, 후자는 붙는다는 것이다. 각각을 사용한 명령어와 실행 모습은 다음과 같다.

\begin{lstlisting}
\[ K + U = \frac{1}{2} mv^2 - \frac{GMm}{r^2}\]
\end{lstlisting}
\[ K + U = \frac{1}{2} mv^2 - \frac{GMm}{r^2}\]

\vskip 2pc

\noindent \begin{lstlisting}
\begin{equation}
S = k_B \ln \Omega \label{equ1}
\end{equation}
\end{lstlisting}
\begin{equation}
	S = k_B \ln \Omega \label{equ1}
\end{equation}

\noindent \textbackslash equation을 사용해서 작성한 식을 언급하기 위해서는 \textbackslash ref\{equ1\}과 같이 ref 명령어를 사용한다. \\
예시 : 통계역학적 엔트로피는 식 \ref{equ1}(\textbackslash ref\{equ1\})과 같이 표현된다.

\vskip 2pc

\noindent 수식이 길어서 한 줄에 표현할 수 없는 경우에는 align 명령어를 사용한다.
\begin{lstlisting}
\begin{align}
F &= \frac{dp}{dt} \notag\\
&= \frac{d}{dt}(mv) \notag \\
&= \frac{dm}{dt}a + m\frac{dv}{dt} \notag \\
&= ma \textrm{ (if mass is constant)} \label{equ2}
\end{align}
\end{lstlisting}
위 명령어를 입력한 후 컴파일하면

\begin{align}
F &= \frac{dp}{dt} \notag\\
&= \frac{d}{dt}(mv) \notag \\
&= \frac{dm}{dt}a + m\frac{dv}{dt} \notag \\
&= ma \textrm{ (if mass is constant)} \label{equ2}
\end{align}

\noindent 위와 같이 여러 줄로 이루어진 수식 \ref{equ2}가 나온다. \\
등호 앞에 있는 \& 기호는 \&를 기준으로 줄을 맞추라는 의미이며 \\
중간 과정의 수식에는 번호를 붙이지 않게 하기 위해 마지막 줄을 제외한 각 줄의 끝에 \textbackslash notag를 붙인다.\\ % 이론적 배경
	\section{연구 과정}
%\section{Method of the study}

\subsection{제목}

내용 % 연구 과정
	\section{결과 및 토의}
%\section{Result and discussion}

\subsection{제목}

내용 % 결과 및 토의
	\section{결론}
%\section{Conclusion}

\subsection{제목}

내용 % 결론
	% \section{추후 연구}
%\section{Further research}

\subsection{제목}

내용 % 추후 연구
	\section{부록}
%\section{Appendix} % 부록
	
	\bibliography{bibfile} % 참고문헌
	% BibTeX 코드 쉽게 얻어오는 방법 %
	% Google Scholar 에서 검색한 결과에서 `인용'을 클릭한다.
	% BibTeX 코드를 얻고자 한다면, 하단의 `BibTeX' 을 클릭.
	% 코드가 나온다. Ctrl+A, Ctrl+C로 복사, bibfile에 붙여넣기.


\end{document}
